\documentclass[11pt]{article}
\usepackage[T1]{fontenc}
\usepackage{lmodern,microtype,amsmath,amssymb,amsthm,mathtools,booktabs,tabularx,array}
\usepackage[letterpaper,margin=0.91in,top=0.78in,bottom=0.83in]{geometry}
\usepackage[dvipsnames]{xcolor}
\usepackage{hyperref,fancyhdr,enumitem}
\definecolor{ink}{HTML}{183747}
\definecolor{teal}{HTML}{126867}
\hypersetup{colorlinks=true,linkcolor=teal,urlcolor=teal,citecolor=teal,pdftitle={Three simplifications of the theta-lift proof of RH},pdfauthor={Research note prepared for Vladimir Reshetnikov}}
\setlength{\headheight}{14pt}
\pagestyle{fancy}\fancyhf{}\fancyhead[L]{\small\color{ink}Three simplifications of the theta-lift proof}
\fancyhead[R]{\small\color{ink}RH and Dirichlet GRH}\fancyfoot[C]{\small\thepage}
\renewcommand{\headrulewidth}{0.25pt}
\setlength{\parindent}{1em}\setlength{\parskip}{2.5pt plus 0.5pt}
\setlist{nosep,leftmargin=*}
\setcounter{tocdepth}{1}
\newtheorem{theorem}{Theorem}[section]
\newtheorem{lemma}[theorem]{Lemma}
\newtheorem{proposition}[theorem]{Proposition}
\newcommand{\NN}{\mathbb N}
\newcommand{\OO}{\mathcal O}
\newcommand{\vv}{\operatorname{v}}
\newcommand{\ord}{\operatorname{ord}}
\newcommand{\HH}{\mathop{\rm Re}}
\newcommand{\pin}{d8415e76ad850651f8c0083afe689dd9f4a45a10}
\newcommand{\source}[3]{\href{https://github.com/VladimirReshetnikov/ProveIt/blob/d8415e76ad850651f8c0083afe689dd9f4a45a10/Analysis/GeneralizedRiemannHypothesis/vendor/grh-lean/grh/src/#1.lean\#L#2}{#3}}
\begin{document}
\thispagestyle{plain}
{\sffamily\footnotesize\bfseries\color{teal} RESEARCH NOTE\hfill 27 NOVEMBER 2026}\par
\vspace{6pt}
{\fontsize{24}{29}\selectfont\bfseries\color{ink} Three simplifications of the\\theta-lift proof of RH\par}
\vspace{5pt}
{\large A sharp dimension estimate, torsion interpolation, and a quadratic sieve}\par
\vspace{10pt}
\noindent\textbf{Abstract.}
The Alp\"oge theta-lift proof compares the volumes of two descriptions of a family of automorphic forms. We give three replacements within this argument. First, at the same allowance for translates as the source's Theorem~B, the unspanned fraction of the target is at most
\[
             19b^{-4}+b^{-5}+2b^{-8}
\]
in place of $4b^{-2}$. A cyclic representation attains the leading constant $19$ asymptotically under the stated moment hypotheses. Second, finite Lagrange interpolation on a sparse set of cyclotomic points gives the simultaneous torsion-specialization theorem needed to control denominators. Its order bound is independent of both the degrees and the number of tests. Third, Heath--Brown's quadratic large sieve and a short M\"obius sum select the auxiliary prime that exposes a proposed zero. We carry the replacements to the Riemann hypothesis and indicate their scope for Dirichlet GRH. The new arguments are proved on paper, with a standalone Lean check of the scalar dimension recurrence and the extremal example's arithmetic. We use the original automorphic estimates as inputs.
\vspace{2pt}
\tableofcontents
\vspace{3pt}
\noindent\textit{Source edition.} We use the readable package in ProveIt, at the snapshot used by the attached 79-page guide \cite{guide,archive}. It corresponds to \href{https://github.com/lsernau/grh-lean/tree/72fead102f663d5c7a407964a13dc0cf9d977207/grh}{upstream revision \texttt{72fead10}}. The verification record identifies this exact edition \cite{check}. The enclosed source index gives declaration-level anchors.
\newpage
\section{Where the three arguments enter}
Fix a multiple $Q$ of $18480$ and a prime $d\equiv Q/2-1\pmod Q$. Write $\eta_d(n)=(n/d)$ and $h_Q=\varphi(Q)$. With $\Lambda$ the von Mangoldt function, define
\begin{equation}
 R_{u,d}(x)=\sum_{\substack{n\le x\\n\equiv1\bmod u}}
 \Lambda(n)(1-\eta_d(n))\left(\frac1n-\frac1x\right)
      -\frac{\log x}{\varphi(u)}-c_{u,d},\qquad u\in\{1,Q\}.
 \label{eq:remainder}
\end{equation}
The fixed-modulus prime number theorem determines $c_{u,d}$ by $R_{u,d}(x)\to0$. The weight $1-\eta_d$ retains the odd powers of inert primes with coefficient two, cancels the split primes, and assigns coefficient one at the ramified prime. In the theta construction these prime powers enter through divisibility of integral derivatives. Two distinct derivative orders on the set of embeddings produce the two cutoffs in the identity
\[
  \mathcal R/a=|B|\bigl(R_{Q,d}(x)-R_{Q,d}(x_1)\bigr)+|A|R_{1,d}(x_1).
\]
Here $a$ is the long derivative order, $A,B$ are two specified sets of embeddings, $x$ is the observation scale, and $x_1$ is a much more distant cutoff. The $m$ successive lifts detect $m\mathcal R$, with $m$ of order $\sqrt{x}$.
\par\noindent
\begin{tabularx}{\linewidth}{@{}p{.20\linewidth}X@{}}
\toprule
\textbf{Part of the proof}&\textbf{Contribution to the prime estimate}\\\midrule
Short forms (A)&An upper bound for the lengths of normalized lifts contains $-m\mathcal R$. The normalization records divisibility and denominators.\\
Enough forms (B)&Few group translates of these vectors span almost all target dimensions. Section~\ref{sec:spectral-spanning} sharpens this assertion, including its translate budget.\\
Long forms (C)&The Rallis inner-product formula, the control of congruences, and a product formula bound the volume from below. Section~\ref{sec:torsion} gives a new proof of one specialization lemma used to select local characters for this comparison.\\
Detection of zeros&The determinant bounds $R_{Q,d}$ from above. Section~\ref{sec:analyticsieve} chooses $d$ so that the resulting transform has an uncancelled pole at the proposed zero.\\\bottomrule
\end{tabularx}
\par
The characters in the last two rows are different. The local characters in the determinant must work through every step of a fixed chain at one coefficient prime. The quadratic character in the prime sum varies with the complex zero under consideration. Confusing the quantifiers would make either replacement insufficient.
\par
For the usual RH the arithmetic construction permits the single modulus $Q=18480$. The principal character at this modulus detects every zeta zero in the open critical strip. This specialization still involves all $\varphi(18480)=3840$ characters through the orthogonality relation for $1\bmod Q$. For Dirichlet GRH the same arguments take $Q=18480q$ when testing a character modulo $q$.
\section{A direct rank estimate on the ramified source}
\label{sec:spectral-spanning}

The short integral forms lie in a specified subspace of the source.
Each new translate of that subspace increases the allowance for denominators.
The question is how much of the target can be filled on the translate budget
of Theorem~B. Keeping the original fourth-power error estimates until the
end gives an upper bound of $(19+o(1))b^{-4}$ for the missed fraction on that budget.
The statement in \texttt{ShortAndEnough.lean} uses $4b^{-2}$.
The coefficient $19$ is optimal for the two moment inequalities and source
codimension estimate used here.

\paragraph{Inputs.}
Let $S,H$ be finite-dimensional complex Hilbert spaces, $h=\dim S>0$,
and let $G$ act unitarily on $H$ by $U_s$. An orthogonal projection $E$
commutes with every $U_s$. Put $V=EH$, $d=\dim V>0$, and let $P:S\to H$
denote the scaled Poincar\'e map. In the source coordinates let $S_r\subset S$
be the space of ramified forms and put $h_r=\dim S_r$. Write
$\mathbb E_s$ for the uniform average over $G$ and $\Pi_L$ for orthogonal
projection onto $L$. The three inputs are
\begin{align}
 \|EP\Pi_L\|_{\rm HS}^2&\ge\dim L-\alpha h
      &&\text{for every subspace }L\subset S,\label{sp:near}\\
 d\,\mathbb E_s\|P^*U_sP\|_{\rm HS}^2&\le(1+\beta)h^2,\label{sp:mom}\\
 h_r&\ge(1-\gamma)h.\label{sp:ram}
\end{align}
The bounds on the first two lines are the source's \source{New/Spanning/SecondMoment/SpanningTheorem}{257}{near-isometry} and \source{New/Spanning/SecondMoment/SpanningTheorem}{288}{second-moment estimate}. The \source{New/Reduction/ShortAndEnough}{410}{ramified-source argument} gives the third. The values are $\alpha=5b^{-4}$, $\beta=9b^{-4}$ and $\gamma=b^{-8}$.
Applying \eqref{sp:near} to $\ker(EP)$ gives
$d\ge(1-\alpha)h$, hence $h\le2d$ for $b\ge3$.

\paragraph{A single recurrence.}
Suppose the selected images span $F\subset V$, with missing fraction
$q=1-\dim F/d$. For $A_s=EU_sP|_{S_r}$ put $L_s=A_s^{-1}F$.
Adding $\operatorname{im}A_s$ gains exactly $g_s=h_r-\dim L_s$
dimensions, by rank--nullity modulo $F$. On $L_s$, the map $EU_sP$
agrees with $\Pi_F U_sP$ because its values lie in $F$ and the component
discarded by $E$ is perpendicular to $F$. Unitarity and commutation
therefore give
\[
 h_r-g_s-\alpha h
 \le \|EU_sP|_{L_s}\|_{\rm HS}^2
 \le \|\Pi_F U_sP\|_{\rm HS}^2 .
\]
This use of the actual source $S_r$ charges its missing dimensions once
in the growth law.

Let $D=\mathbb E_s U_sPP^*U_s^*$. Expanding the square and grouping
pairs of elements by $s^{-1}t$ gives
$\operatorname{tr}(D^2)=\mathbb E_s\|P^*U_sP\|_{\rm HS}^2$.
The average of the last squared norm in the preceding display is
$\operatorname{tr}(\Pi_FD)$. Trace Cauchy--Schwarz and \eqref{sp:mom}
bound this quantity by $h\sqrt{(1-q)(1+\beta)}$. Consequently
\begin{align}
 \mathbb E_s g_s
  &\ge h\{1-\alpha-\gamma-\sqrt{(1-q)(1+\beta)}\}\label{sp:exactgrowth}\\
  &\ge \tfrac h2(q-\nu),\qquad
    \nu=\beta+2\alpha+2\gamma=19b^{-4}+2b^{-8}.\nonumber
\end{align}
The last line uses $2\sqrt u\le1+u$ and $q\beta\ge0$.
Choose a translate with at least the average gain. The new fraction obeys
$q_{j+1}\le(1-h/(2d))q_j+(h/(2d))\nu$. Starting from $q_0=1$
and using $0<h/(2d)\le1$ gives, after $n$ choices,
\begin{equation}
 q_n\le\nu+(1-h/(2d))^n\le\nu+\exp(-nh/(2d)).
 \label{sp:decay}
\end{equation}
Repeated translates can be discarded without changing the span.

\paragraph{A stronger rank bound on the same budget.}
Take $n=\lceil(10d/h)\log b\rceil$ with $b\ge3$. For the resulting
set $J$ the cost and missing fraction satisfy
\begin{align}
 |J|h_r &\le nh\le 10d\log b+2d\le5d\log(1+2b^2),
   \label{sp:card}\\
 q_n&\le 19b^{-4}+b^{-5}+2b^{-8}<20b^{-4}.
   \label{sp:rank}
\end{align}
For the last inequality in \eqref{sp:card} use $5\log 2>2$.
Equations \eqref{sp:card}--\eqref{sp:rank} imply the original two
allowances, since $20b^{-4}\le4b^{-2}$ for $b\ge3$.
The coordinate isomorphism and the nonzero scalar normalizations identify
these images of $S_r$ with the images used by Theorem~B. Independent
vectors can still be selected from the actual translates of the short
integral forms supplied by Theorem~A.

\paragraph{The sharp obstruction from the two moments.}
Take $F$ to be the sum of \emph{all} allowed images. Then $g_s=0$ for
every $s$, so \eqref{sp:exactgrowth} gives a permanent missed fraction of
at most
\[
 q_* = 1-\frac{(1-\alpha-\gamma)^2}{1+\beta}
 \qquad(\alpha+\gamma<1).
\]
For full source this threshold can hold with equality. Let $b\ge3$ be
an integer, $x=b^4$, $m=(x-5)^2$, and $d=x(x+9)$. On $V=H=\mathbb C^d$,
a cyclic group of order $m$ permutes the first $m$ coordinate axes
regularly and fixes the other $d-m$. Take $E=I$, $S_r=S=\mathbb C$
and $P(1)=\sqrt{1-5/x}\,e_1$. The near-isometry condition is sharp.
Only the identity group element has $P^*U_sP\ne0$, so the normalized
second moment is
\[
 d\,\mathbb E_s\|P^*U_sP\|_{\rm HS}^2
       =\frac{d(1-5/x)^2}{m}=1+9/x.
\]
The whole orbit spans exactly the first $m$ axes. Its missed fraction is
\[
 \frac{19x-25}{x(x+9)}
       =1-\frac{(1-5/x)^2}{1+9/x}
       =(19+o(1))b^{-4}.
\]
Thus any universal leading coefficient smaller than $19$ fails for these
three input inequalities, even with access to every translate. Further
arithmetic information about the Poincar\'e map could produce a stronger
conclusion in the actual automorphic spaces. The example determines the
limit of the moment argument. For a concrete comparison of the abstract estimates, at $b=10$ the published allowance is $4\%$, the new upper bound is $0.191002\%$, and this cyclic example actually misses about $0.189804\%$.

\paragraph{Relation to the existing proof.}
The group average, Cauchy--Schwarz bound, and preimage rank identity are
already present in \texttt{FewTranslates.lean}. The change is to apply
them on $S_r$ in the first step, retain the original numerical errors,
and solve one affine recurrence. The existing proof first weakens both
moment errors to $(8b^2)^{-1}$, obtains enough translates of the entire
source, and then bounds the cumulative loss on restricting their images.
The local orbital estimates enter at precisely \eqref{sp:near} and \eqref{sp:mom}. Thus a formal refactor can reuse the source's \source{New/Spanning/SecondMoment/FewTranslates}{288}{Hilbert--Schmidt estimate} and its \source{New/Spanning/SecondMoment/FewTranslates}{393}{rank identity}, then change the dimension recurrence and the final restriction step.

\section{Simultaneous torsion specialization by interpolation}
\label{sec:torsion}
In the determinant lower bound one needs a single local character whose first $m$ powers keep finitely many valuations small. Both special values and Selmer relations give polynomials in the character coordinates. Their degrees can depend on the level at which the character is constructed. The bound on the level must be chosen first. The source handles this by factoring a polynomial, counting its roots on valuation circles, and estimating their truncated distances to primitive roots of unity. We give a proof by interpolation on a set chosen before seeing any polynomial.

Fix a prime $\ell$. Work in $\overline{\mathbb Q}_\ell$ with its ordinary additive valuation $v(\ell)=1$, $v(0)=+\infty$, and valuation ring $\OO$. For $h\ge1$ let $t_h=1/\varphi(\ell^h)$. Every primitive $\ell^h$-th root $\xi$ satisfies $v(\xi-1)=t_h$, since its conjugates have equal valuation and $\Phi_{\ell^h}(1)=\ell$.

\begin{theorem}[Uniform torsion specialization]\label{thm:torsion}
Fix $s_0\in\NN$, $r\in\mathbb Q_{>0}$, $B\ge0$ and $\varepsilon>0$. There is an $H$ with the following property. Let $n\le s_0$ and let $P_i\in\OO[X_1,\ldots,X_n]$ be any finite family. Suppose $x_i\in\OO^n$, all $v(x_{ij})\ge r$, and $\sum_i v(P_i(x_i))\le B$. There exist $\ell^H$-th roots of unity $\zeta_1,\ldots,\zeta_n$, with $v(\zeta_j-1)<r$, such that
\[
 v\bigl(P_i(\zeta_1-1,\ldots,\zeta_n-1)\bigr)
                         \le v(P_i(x_i))+\varepsilon\qquad\text{for every }i .
\]
\end{theorem}

\subsection*{A finite set on which interpolation costs little}
Choose $0<\delta\le1$. For every $h$ with $t_h<r$, fix a primitive root $\xi_h$ of order $\ell^h$. At this level use the points
\[
 z=\xi_h^{a}-1,\qquad 1\le a\le N_h,\quad \ell\nmid a,
 \qquad N_h=\left\lfloor\frac{\delta}{4(h+1)t_h}\right\rfloor.
\]
Each point has weight $t_h$. Write $\mu(S)=\sum_{z\in S}v(z)$ for the weight of a finite set of nodes. Nodes from distinct levels have distinct valuations. Within one level $N_h<\ell^h$, so distinct exponents yield distinct nodes. The weight available over all levels is infinite. In fact $N_h$ tends to infinity, and at least $N_h/2$ of the first $N_h$ integers are prime to $\ell$. For all large $h$ the weight at level $h$ is at least $\delta/(16(h+1))$. Thus for any prescribed mass we can fix a finite union of levels exceeding it.

For two different nodes at level $h$ we have
\[
 v(\xi_h^a-\xi_h^{a'})=t_h\ell^{v_\ell(a-a')}.
\]
For a fixed $a$, the sum of these valuations over the other nodes is at most
\begin{equation}
 2t_h\sum_{k=1}^{N_h}\ell^{v_\ell(k)}
 \le 2t_hN_h(1+\lfloor\log_\ell N_h\rfloor)
 \le\delta/2. \label{eq:spacing}
\end{equation}
The middle bound follows by summing $\ell^j\lfloor N_h/\ell^j\rfloor$ over the possible $j$. Only levels with nodes enter this calculation.

\begin{lemma}[Weight of unsuccessful nodes]\label{lem:capacity}
Let $g\in\OO[X]$, $v(y)\ge r$, and $b=v(g(y))<\infty$. In any finite union of the above levels, the nodes with $v(g(z))>b+\delta$ have total weight at most $b$.
\end{lemma}
\begin{proof}
Take the set $S$ of these nodes and suppose $\mu(S)>b$. Divide $g$ by the monic polynomial $D(X)=\prod_{w\in S}(X-w)$. Its quotient $q$ and remainder have integral coefficients. Lagrange interpolation of the remainder gives the exact identity
\[
 g(y)=q(y)D(y)+\sum_{z\in S}g(z)
            \prod_{w\in S\setminus\{z\}}\frac{y-w}{z-w}.
\]
Since $v(y)>v(w)$, $v(D(y))=\mu(S)>b$. In each interpolation factor $v(y-w)=v(w)$. For nodes of different levels $v(z-w)=\min(v(z),v(w))$, so their contributions to the valuation of the product are nonnegative. The total loss at the same level is at most $\delta/2$ by \eqref{eq:spacing}. Every interpolation summand and the quotient term have valuation greater than $b$. The ultrametric inequality contradicts $v(g(y))=b$.
\end{proof}

The bound on the weight is exact. For any subset $S$ take $g_S(T)=\prod_{w\in S}(T-w)$ and $y=0$. Then $v(g_S(0))=\mu(S)$ and every point of $S$ is unsuccessful. As each grid weight is positive, Lemma~\ref{lem:capacity} says the unsuccessful set is exactly $S$. In particular a fixed grid of weight at most $B$ can be defeated even by one test with reference valuation at most $B$. This explains the strict mass threshold in the selection argument.

\subsection*{Moving many different reference points at once}
Suppose the $g_i$ are integral one-variable polynomials with $v(y_i)\ge r$, $b_i=v(g_i(y_i))$ and $b_*=\sum_i b_i<\infty$. One extra test controls their total valuation after the move. Choose $u$ with $v(u)=r$. For each $g_i(uT)$ divide by a coefficient of least valuation. Its reduction is a nonzero polynomial over $\overline{\mathbb F}_\ell$. There is a single nonzero residue class avoiding the roots of all these reductions. Lift it to a unit $w$. At $y_*=uw$, every $g_i$ attains this least coefficient valuation. As $y_i/u$ is integral, $v(g_i(y_*))\le b_i$. The product $G=\prod_i g_i$ therefore satisfies $v(G(y_*))\le b_*$.

Apply Lemma~\ref{lem:capacity} to each $(g_i,y_i)$ and then to $(G,y_*)$. The union of all unsuccessful nodes has weight at most $2b_*$. On any fixed grid of weight greater than $2b_*$, some node $z$ therefore satisfies both
\[
 v(g_i(z))\le b_i+\delta\quad(\text{all }i),\qquad
 \sum_i v(g_i(z))=v(G(z))\le b_*+\delta.
\]
The second inequality prevents the number of tests from entering the level bound.

To prove Theorem~\ref{thm:torsion}, take $\delta=\min(1,\varepsilon)/(s_0+1)$ and fix a grid of weight greater than $2(B+1)$. Its greatest level is the promised $H$. Freeze every coordinate except the next one at the current point of each test. Apply the one-variable conclusion and update this coordinate to the same node in every test. Unvisited coordinates still have valuation at least $r$. After $j$ moves the sum is at most $B+j\delta\le B+1$, and each individual increase is at most $j\delta\le\varepsilon$. This proves the assertion even as the family size and polynomial degrees vary.

\subsection*{Every power of the same character}
For each $1\le k\le m$, apply the theorem simultaneously to the polynomials
\[
 P_{i,k}(Y_1,\ldots,Y_n)=P_i\big((1+Y_1)^k-1,\ldots,(1+Y_n)^k-1\big).
\]
For an old coordinate $x$ pick a root $y$ of $(1+Y)^k-(1+x)$ satisfying $v(y)\ge r/k$. If $x=0$, take $y=0$. Otherwise the polynomial is monic and integral with constant $-x$, so its $k$ roots are integral and the sum of their valuations is $v(x)$. At least one has the required valuation. Use radius $r/m$ and total budget $mB$. We obtain one tuple with $v(\zeta_j-1)<r/m$ whose $k$-th powers satisfy the requested inequalities for every $i,k$. In the application $r<1$. An order at most $m$ would give $v(\zeta_j-1)\ge 1/m$, so each power is nontrivial. The exact formula for the order of $\zeta_j^k$ gives $v(\zeta_j^k-1)=\ell^{v_\ell(k)}v(\zeta_j-1)\le r$.

\noindent\textit{Source interface.} This proves the conclusion of \source{New/Congruences/Selmer/RootsOfUnity}{474}{the multivariable torsion theorem} and supplies the hypothesis used in \source{New/Congruences/Selmer/GeneralTheorem}{325}{the all-powers construction}. Its bound $H$ must precede the auxiliary level, because the source constructs the polynomials at that level. Harmonic divergence gives such an $H$ with no useful computational size guarantee. In a formal implementation the Lagrange quotients require the signed valuation on the fraction field, or an equivalent norm inequality after clearing denominators. All evaluations in the theorem's conclusion are integral.

\section{A finite Dirichlet polynomial selects the auxiliary prime}
\label{sec:analyticsieve}
We give a replacement for the quadratic-twist selection theorem used at the end of the argument. It applies to every Dirichlet character of the fixed modulus, including the complex characters. Its substantial external input is Heath--Brown's quadratic large sieve. The proof only needs the following form \cite[Corollary~2]{HB}. If $\mathcal{S}(U)$ consists of the primitive real characters of conductor at most $U$, then for complex coefficients $a_n$ and every $e>0$,
\begin{equation}\label{eq:hbsieve}
 \sum_{\xi\in\mathcal{S}(U)}\left|\sum_{n\le V}a_n\xi(n)\right|^2
 \ll_e (UV)^e(U+V)\!\sum_{\substack{j,k\le V\\jk\text{ a square}}}|a_j a_k|.
\end{equation}
We average over the real quadratic character $\eta_d$. The fixed, possibly complex, character $\psi$ goes into the coefficients.

\paragraph{Proposition.} Fix $Q$, $1/2<\beta<1$, and $H\ge0$. Among the primes $d\in(X,2X]$ coprime to $Q$, the number for which some $L(s,\psi\eta_d)$, $\psi\bmod Q$, has a zero in $K=[\beta,1]+i[-H,H]$ is $O_{Q,\beta,H}(X^{1-c})$ for some $c=c(\beta)>0$. Here $\eta_d(n)=(n/d)$ for odd prime $d$. For example, $c=(\beta-1/2)^2/8$ is admissible. Consequently every invertible residue class modulo $Q$ contains good primes, including $Q/2-1\bmod Q$ for $18480\mid Q$.

\paragraph{Proof.} Write $\alpha=(\beta-1/2)/2$, $\sigma_0=1/2+\alpha$ and $b=\alpha/2$. Fix a bounded open rectangle $U_0$ containing $K$ with $\Re s>\sigma_0$ throughout. It contains a disk of some fixed radius $r>0$ about each point of $K$. Fix $\psi\bmod Q$ and write $\chi=\psi\eta_d$. Its partial sums on every interval are $O_Q(\sqrt X\log(2X))$. Indeed partition the interval by classes $n=j\bmod Q$. In each class $\psi$ is constant and $\eta_d(Qk+j)=\eta_d(Q)\eta_d(k+jQ^{-1})$, with $Q^{-1}$ taken modulo $d$. The classical P\'olya--Vinogradov inequality applies to each of the resulting intervals. Partial summation therefore gives, uniformly on $U_0$,
\[
 L(s,\chi)=\sum_{n\le N}\frac{\chi(n)}{n^s}
     +O_{Q,U_0}\big(\sqrt X\log(2X)N^{-\sigma_0}\big).
\]
This follows from the ordinary Dirichlet series on $\Re s>0$ for a nonprincipal character. These twists are nonprincipal since $d\nmid Q$ and the quadratic character at $d$ is nonprincipal.

Take $N=\lfloor X^{1-b}\rfloor$, $M=\lfloor X^b\rfloor$ and $P_d(s)=\sum_{m\le M}\mu(m)\chi(m)m^{-s}$. Since $|P_d(s)|\ll M^{1-\sigma_0}$, putting $F_d=LP_d-1$ gives
\begin{align}
 F_d(s)&=\sum_{M<k\le NM} A(k)\psi(k)k^{-s}\eta_d(k)
                 +O_{Q,U_0}\big(X^{-\alpha/2}\log(2X)\big),\label{eq:shortproduct}\\
 A(k)&=\sum_{\substack{m\mid k\\m\le M,\ k/m\le N}}\mu(m).\nonumber
\end{align}
Indeed $A(1)=1$, $A(k)=0$ for $1<k\le M$ by M\"obius inversion, and $|A(k)|\le\tau(k)$. The exponent in the error is $1/2-(1-b)\sigma_0+b(1-\sigma_0)=-\alpha+b$. The two truncations have complementary lengths, so $NM\le X$.

For $a_k=\mathbf 1_{M<k\le NM}A(k)\psi(k)k^{-s}$ the square-pair sum in \eqref{eq:hbsieve} satisfies, for any $0<e<\alpha$,
\begin{align*}
 \sum_{jk=\square}|a_j a_k|
 &\ll_e\sum_{h\ \mathrm{squarefree}}
   \left(\sum_{h v^2>M}(h v^2)^{-\sigma_0+e}\right)^2
 \ll_{e,\sigma_0} M^{-2\alpha+2e}\log(2M).
\end{align*}
For $h\le M$, comparison with the integral of $v^{-2\sigma_0+2e}$ bounds the inner sum by $O(M^{-\alpha+e}h^{-1/2})$. For $h>M$ it is $O(h^{-\sigma_0+e})$. The standard divisor estimate $\tau(k)\ll_e k^e$ gives the first inequality. It is essential to count all pairs with square product, since different integers can have the same squarefree part.

Apply \eqref{eq:hbsieve} with $U=2X$ and $V=NM$, putting $\psi$ in the coefficients. Summing the squared uniform error in \eqref{eq:shortproduct} over at most $2X$ primes gives
\[
 \sum_{X<d\le2X\ {\rm prime},\ d\nmid Q}|F_d(s)|^2
  \ll_{Q,U_0,e} X^{1-2b\alpha+2e+2be}\log(2X)
                +X^{1-\alpha}\log^2(2X).
\]
Here and above $X$ can be taken sufficiently large for $N,M\ge1$. Since $2b\alpha=\alpha^2$, take $e$ sufficiently small in terms of $\alpha$ and absorb logarithms to obtain $O(X^{1-\alpha^2/2})$, uniformly in $s\in U_0$.

If $L(z,\chi)=0$ for some $z\in K$, then $F_d(z)=-1$. The mean-value inequality for a holomorphic function gives
\[
 1\le (\pi r^2)^{-1}\int_{U_0}|F_d(s)|^2\,dA(s).
\]
The radius and the containing rectangle work uniformly as the zero varies with $d$. Integrating the uniform bound above counts at most $O(X^{1-\alpha^2/2})$ such primes. A union over the finitely many $\psi\bmod Q$ proves the bound with $c=\alpha^2/2$. The fixed-modulus prime number theorem gives $\asymp_Q X/\log X$ primes per invertible progression in $(X,2X]$. This dominates the exceptional count. For $\Re s>1$ all twists are nonzero by the Euler product.\hfill$\square$

One can also finish with the exact prime-supply input used in the source. On $(2^j,2^{j+1}]$ the exceptional primes contribute $O(j2^{-cj})$ to $\sum_d (\log d)/d$. This is summable. The same weighted series diverges over any invertible progression, so it contains arbitrarily large good primes. If the requested box begins at $\beta\ge1$, a box with left edge strictly between $1/2$ and $1$ together with the Euler product suffices.

\noindent\textit{Source interface.} The proposition supplies the conclusion of \source{Lib/Analytic/QuadraticTwist}{796}{the choice of a good prime}. The source proves a stronger quantitative exponent in its parametrization. Only existence is consumed by \source{New/Reduction/Zeros/Criterion}{37}{the one-sided criterion}. The large sieve replaces the auxiliary Gaussian theta-series estimate and strip interpolation in this branch. Its own proof is a substantial external input. The reference is Corollary~2 of \cite{HB}, which allows arbitrary complex coefficients. Invoking a zero-density theorem restricted to real target characters would leave a gap for complex $\psi$.

\section{Assembly and the Riemann hypothesis}\label{sec:endpoint}
\subsection*{What the determinant proves}
In the notation of the guide let $d_T$ be the dimension of the target, $f$ the number of infinite places in the chosen CM type, $a$ the derivative order, and $b$ the larger ordinary weight. Set $L=fa(\log b)^4$, $I=10\log(1+2b^2)fa(\log b)^3$. Let $H_0\ge0$ denote the common volume main term. Theorem~A gives sufficiently short vectors. Section~\ref{sec:spectral-spanning} permits a choice of $n$ independent vectors from their translates at the allowance required for Theorem~C. The local character construction in Section~\ref{sec:torsion}, together with the source's cohomological estimates and special-value identities, supplies the requisite denominator bound for C. The squared-volume comparison is
\[
 d_T(H_0-3L)\le n(H_0-m\mathcal R+L)+2d_T I.
\]
The determinant on the upper side is bounded by the product of squared lengths. On the lower side the $2d_T I$ term accounts for clearing denominators. Both sides carry the same $H_0$, a mathematical equality furnished by the one-lift period and special-value formulas. Since $n\le d_T$ and $n\ge d_T/2$ for the large scales at issue,
\[
 nm\mathcal R\le(n-d_T)H_0+(n+3d_T)L+2d_T I
              \le d_T(4L+2I).
\]
For the source's large $b$ this gives \source{New/Reduction/MainInequality}{87}{$m\mathcal R\le fa(\log b)^5$}. The tighter rank estimate gives extra margin here. The gain in the power of $x$ still comes from dividing by the number $m$ of lifts.

Substitute the two-cutoff identity from Section~1 and use $f=\varphi(Q)|A|$, $|A|/|B|\le\log b$. Suppressing the fixed $d$ in the subscripts, we get
\[
 R_Q(x)\le\frac{\varphi(Q)(\log b)^6}{m}
                 +|R_Q(x_1)|+(\log b)|R_1(x_1)|.
\]
For any sufficiently small fixed $\tau>0$ the source supplies data at every large integer $x$ with $m\ge\sqrt x$, $\log b\le 2x^\tau$ and $\log x_1\ge x^{\tau/12}$. At the fixed moduli involved, the classical prime number theorem with its fixed-modulus remainder gives $|R_u(y)|\ll_A(\log y)^{-A}$ for every fixed $A$. One can also see the smoothing step directly: for $T_u(t)=\sum_{n\le t,\ n\equiv1\ (u)}\Lambda(n)(1-\eta_d(n))$, integrate $(T_u(t)-t/\varphi(u))/t^2$ from $y$ to infinity and reverse its sign. The constant in \eqref{eq:remainder} makes this tail integral equal to $R_u(y)$. The estimate for $T_u$ follows from the fixed-modulus progression theorem modulo $ud$ with finitely many ramified prime powers handled separately. Taking $A$ large enough absorbs the far terms. Between consecutive integers the smoothed remainder changes by $O(1/x)$ by the Chebyshev bound. Taking $6\tau<\varepsilon$ recovers the arithmetic conclusion, for every fixed admissible $Q,d$ and $\varepsilon>0$:
\begin{equation}
                  R_{Q,d}(x)\le x^{-1/2+\varepsilon}
                     \quad\text{for all sufficiently large real }x.
 \label{eq:finalbound}
\end{equation}
The upper bound with its $d$-dependent threshold suffices for zero detection.

\subsection*{Why the new prime selection suffices}
Suppose that a character at modulus $Q$ has a zero $\rho$ in $1/2<\Re\rho<1$. Choose $1/2<\beta<\Re\rho$. Section~\ref{sec:analyticsieve} supplies a prime $d\equiv Q/2-1\pmod Q$ such that every twisted function is nonzero for $\Re s\ge\beta$, $|\Im s|\le |\Im\rho|$. Fix this prime and apply \eqref{eq:finalbound}. This order of choices is essential.

Let $E_\psi=L(s,\psi)$ for a nonprincipal character, and remove the pole of a principal character by replacing its $L$-function with $(s-1)L(s,\psi)$ regularized at one. Let $n(z)$ and $n_d(z)$ count the zeros at $z$, with multiplicity, summed over all $E_\psi$ and $E_{\psi\eta_d}$ respectively. Orthogonality in the class $1\bmod Q$ gives every character coefficient one. The smoothed transform
\[
 \Phi(s)=\int_0^\infty R_{Q,d}(e^v)e^{(1-s)v}\,dv
\]
initially exists far to the right. The integral of a single kernel is $n^{-s}/[s(s-1)]$. Insert the Dirichlet series for the logarithmic derivatives and continue meromorphically. At $z\ne0,1$ in $\Re z>0$ the residue is
\[
             \frac{n(z)-n_d(z)}{\varphi(Q) z(1-z)}.
\]
The centering makes the singularity at one removable. These identities are proved in the source's \source{New/Reduction/Zeros/PrimeError}{414}{prime-error module}. The chosen $d$ makes the residue nonzero at $\rho$. It also makes it nonnegative at every real point $u\in(\beta,1)$. Regularized functions have no zeros on the real interval $[1,\infty)$.

For completeness take $K$ so large that $K e^{\beta v}-e^vR_{Q,d}(e^v)\ge0$ on $[0,\infty)$, using \eqref{eq:finalbound} with $\varepsilon=\beta-1/2$ and continuity on a bounded initial interval. The Laplace transform of this positive density continues as $K/(s-\beta)-\Phi(s)$. At a finite real abscissa of convergence $a>\beta$, its only possible singularity is a simple pole of nonpositive residue, by the preceding sign computation. A nonnegative transform has nonnegative residue there, by taking the real limit from the right. Thus such a singularity would be removable. Landau's lemma excludes removability at a finite abscissa of convergence of a positive measure: a Taylor expansion just to its right has nonnegative integrated terms when evaluated to the left, and convergence across the abscissa would extend the integral there by monotone convergence. Hence the integral is holomorphic throughout $\Re s>\beta$. The continued $\Phi$ is holomorphic at $\rho$, contradicting its nonzero residue.

\subsection*{Specialization and scope}
At $Q=18480$ the admissible residue class is $9239\pmod{18480}$. The principal character has
\[
 L(s,\mathbf1_Q)=\zeta(s)\prod_{p\mid Q}(1-p^{-s}),
\]
whose finite factors are nonzero for $\Re s>0$. We have excluded every zeta zero in the right half of the open critical strip. The functional equation reflects any zero in its left half into the right half. Together with the classical zero-free boundary, Euler product and trivial-zero classification, this proves RH from the retained inputs. For a character modulo $q$ one inflates to $Q=18480q$. Finite Euler factors preserve its zeros in the open strip, and the functional equation for its primitive representative reflects a zero to one of the inverse character. The argument therefore gives the Dirichlet GRH endpoint as well.

\section{Verification and remaining formal work}
The attached guide's sections 9, 15, 16 and 17 identify the points in the original construction corresponding respectively to Sections~\ref{sec:spectral-spanning}, \ref{sec:torsion}, the volume comparison, and Section~\ref{sec:analyticsieve}. We compared the declarations and their uses with the pinned source. Andrew Sutherland's public record reports a successful build, replay of all 803 readable modules, independent checking of the exported theorems, and comparison with the specified statements \cite{check}. Those checks concern the original edition. We have given mathematical proofs of the replacements. The included \texttt{RankRecurrence.lean} was checked with Lean 4.32.0 and its standard library. It proves the scalar induction over linearly ordered commutative rings and the exact rational identity $x(19x-25)/[x(x+9)]=19-196/(x+9)$ for $x>0$. A kernel-checked implementation of all three replacements in the pinned Lean 4.35.0 package remains open.

\begin{center}\small
\begin{tabularx}{\linewidth}{@{}p{.17\linewidth}X@{}}
\toprule
\textbf{Change}&\textbf{What a Lean implementation would need}\\\midrule
Rank estimate&
Reuse the projected-norm inequality and rank--nullity lemmas. Form the preimage inside the ramified source, keep both fourth-power errors, and prove the scalar recurrence. This change has the smallest external dependency cost.\\
Torsion grid&
Prove the spacing sum, monic division and valuation bound for the Lagrange basis on every subset of the grid. The public theorem can retain its current statement, including the quantifier over all finite index types after the level bound.\\
Good prime&
The source already proves P\'olya--Vinogradov. The proof here also requires the quadratic large sieve with the square-pair term for arbitrary complex coefficients. Its formalization cost could exceed the savings in this branch.\\\bottomrule
\end{tabularx}
\end{center}
The archive includes exact finite checks of the rank calculation and extremal family, cyclotomic distances and interpolation bounds, and the sieve coefficient identities. A separate script verifies the listed Lean declaration anchors and Git blob hashes against a checkout. The README gives the reproduction commands. Measuring any change in formal proof size or build time requires an integrated implementation.

\subsection*{Further research questions}
The extremal cyclic action has a large invariant complement of the orbit. Does a particular automorphic feature of the actual Poincar\'e maps rule out this configuration strongly enough to beat $19b^{-4}$? For the local problem, can one raise the weight of an explicit cyclotomic grid at a fixed maximum order while bounding every Lagrange denominator as in \eqref{eq:spacing}? Both questions have numerical targets. Either improvement would sharpen its own step. Changing the final exponent in \eqref{eq:finalbound} would additionally require an improvement in the determinant or chain-length estimates.

\begin{thebibliography}{9}\small
\bibitem{archive} L. Alp\"oge and collaborators, theta-lift proof of RH and Dirichlet GRH. L. Sernau, \emph{grh-lean}, readable edition, \href{https://github.com/lsernau/grh-lean/tree/72fead102f663d5c7a407964a13dc0cf9d977207/grh}{revision \texttt{72fead10}}. \href{https://github.com/VladimirReshetnikov/ProveIt/tree/d8415e76ad850651f8c0083afe689dd9f4a45a10/Analysis/GeneralizedRiemannHypothesis/vendor/grh-lean}{ProveIt snapshot and vendor notes}.
\bibitem{guide} \emph{How theta lifts locate the zeros of Dirichlet $L$-functions}, supplied guide, 79-page edition, 27 November 2026. See especially sections 9, 15--19 and its source map.
\bibitem{check} A. V. Sutherland, \href{https://math.mit.edu/~drew/shortgrh/}{Verification files for the shorter Lean proof of RH and GRH for Dirichlet $L$-functions}, verification record for \texttt{72fead10}.
\bibitem{HB} D. R. Heath-Brown, \href{https://doi.org/10.4064/aa-72-3-235-275}{A mean value estimate for real character sums}, \emph{Acta Arith.}\ \textbf{72} (1995), 235--275, Corollary 2.
\bibitem{note} L. Alp\"oge, \href{https://alpo.ge/remark.pdf}{Announcement and notes on the discovery of the proof}, November 2026.
\end{thebibliography}

\end{document}
